\section{Introduction}
\label{sec:intro}

Over the past two years, the mechanistic interpretability program has
converged on an unexpected empirical regularity. Sparse autoencoders
(SAEs) trained on the residual stream of large transformers recover
dictionaries of features that, individually, behave as discrete
computational primitives \cite{bricken2023monosemanticity,
cunningham2024sparse, templeton2024scaling, gao2024scaling}. The
features are interpretable across many examples; they activate on
recognizable concepts; their ablation produces predictable behavioral
changes. This is the case across model families, model sizes, training
regimes, and tokenizers. The pattern is sufficiently robust that the
\emph{universality hypothesis}---the conjecture that different networks
trained on similar tasks recover similar internal computations---has
been advanced as a working principle of the field
\cite{olah2020zoom, chughtai2023toy, gurnee2024universal}.

A separate line of work has documented the same regularity from the
representational-geometry angle. Bansal et al.~\cite{bansal2021revisiting}
showed that the lower layers of one model can be ``stitched'' to the
upper layers of another with only a thin trainable interface, and that
the stitched model performs nearly as well as either original; this
implies a deep correspondence between the intermediate representations
of independently trained networks. Kornblith et al.~\cite{kornblith2019similarity}
established centered kernel alignment (CKA) as a similarity measure
robust to the symmetries that should not matter for the comparison and
showed that representations across models are far more similar than the
sheer dimensionality of the representation spaces would predict at
chance. Most recently, Huh et al.~\cite{huh2024platonic} formulated the
\emph{Platonic Representation Hypothesis}: that as models scale and
capability increases, their representations converge toward a shared
statistical model of the world.

What has been missing, in our reading, is the mathematical object that
this empirical convergence corresponds to.

\subsection*{The Operator-Theoretic Bridge}

We propose that the SAE feature dictionary is a concrete realization of
the \emph{approximately invariant subspace structure} of the nonlinear
forward operator that constitutes a transformer's inference pass. The
classical theory of invariant subspaces of bounded linear operators on
Hilbert spaces \cite{lomonosov1973invariant, enflo1987nontrivial,
beauzamy1988introduction, aronszajn1954invariant} asks whether every
such operator admits a non-trivial closed invariant subspace---that is,
a subspace $V$ such that $T(V) \subseteq V$. The question is open in
the Hilbert space setting; Lomonosov \cite{lomonosov1973invariant}
established the affirmative for operators commuting with a non-zero
compact operator, and Enflo \cite{enflo1987nontrivial} showed
counter-examples exist for general Banach spaces. The structural
intuition is the older one: when an operator admits an invariant
subspace, its action decomposes into restrictions to simpler
sub-operators, and the geometry of the operator becomes legible through
the geometry of its invariant subspaces.

Transformer forward passes are not bounded linear operators. They
compose linear projections, attention computations, and pointwise
nonlinearities. The classical theorems do not apply directly. We use
the language of invariant subspaces in an \emph{approximate} sense: a
subspace $V$ is approximately invariant under a (possibly nonlinear)
map $f$ on a representative sample if the projection of $f(V)$ back
onto $V$ leaves a residual that is small compared to the operator's
overall action. This is the operational content of the empirical
finding that SAE features are stable across forward passes: they
identify directions in activation space that the forward operator
preserves approximately.

The mechanistic interpretability literature has been doing operator
theory for two years without naming it as such. Our first contribution
is to make the bridge explicit.

\subsection*{The Falsifiable Consequence}

If the bridge holds---if SAE feature dictionaries are bases for
approximately invariant subspaces of the forward operators of trained
transformers---then a sharp empirical consequence follows. Two models
$M_A$ and $M_B$ with different architectures, run on the same
compositional probe set, should admit an orthogonal alignment between
their respective approximately invariant subspaces, with Procrustes
residual significantly below the chance level set by a permutation null
that breaks the semantic correspondence between rows while preserving
each model's internal geometry. Performance comparability between
architectures (the weak universality claim) is consistent with many
internal mechanisms; representational isomorphism is not. The
isomorphism is testable. The test can return a null. We pre-register
both the protocol and the significance criterion before observing data.

This is the \emph{Cross-Operator Subspace Isomorphism} (COSI) program,
of which the present paper is the calibration step.

\subsection*{What This Paper Reports}

We report a representational-alignment pilot in two pre-registered
runs. We are explicit that the pilot does not test the strong
operator-theoretic claim; it tests a weaker but still interesting
sub-claim, with controls and interventions that would close the gap
specified as an explicit research program (\S\ref{sec:requiredwork}).

\textbf{Phase~0:} Phi-4-Reasoning-Plus, a dense pure-attention
reasoning model in the Phi-3 architecture family \cite{abdin2024phi4},
against Qwen3-4B, a dense pure-attention generalist model in the Qwen3
family \cite{qwen2025qwen3}, on $N = 96$ compositional-evaluation
prompts drawn from a single political-economy domain. The observed
Procrustes residuals are $0.61$, $0.81$, $0.81$ at the three
pre-registered layer fractions; the empirical permutation $p$-value
is $< 0.001$ at every layer (the observed residual is below the
$1$st percentile of the $S = 1000$ permutation null distribution).

\textbf{Phase~1:} Phi-4-Reasoning-Plus against Qwen3-Next-80B-A3B-Thinking
\cite{qwen2025qwen3}, a hybrid attention/state-space-model architecture,
on $N = 600$ prompts split equally across mathematical, logical, and
political-economy compositional-evaluation domains. The observed
Procrustes residuals are $0.77$, $0.91$, $0.96$ at the three layer
fractions, with empirical $p < 0.001$ at every layer and at every
per-domain stratification cell (three domains $\times$ three layers,
$n = 200$ per cell, $S = 500$ permutation samples).

\subsection*{What These Results Do and Do Not Show}

These results show that, on matched compositional-evaluation prompts,
the PCA-projected last-token activation matrices of two architecturally
different frontier transformers admit a substantially better orthogonal
alignment than they do under random row permutation. This is a real
cross-model representational signal. We report it in full.

These results do \emph{not} yet show:

\begin{enumerate}
    \item That the aligned PCA subspaces are approximately invariant
        under the forward operator in the sense the operator-theoretic
        framing requires. The paper defines approximate invariance as
        a leakage statistic but does not measure it; PCA top-$k$
        components capture variance, not necessarily invariance.
    \item That the alignment isolates compositional reasoning structure
        from lexical, template, syntactic, or tokenizer-induced
        structure. The row-permutation null is too weak to discriminate
        these. Two competent language models given the same English
        prompts will produce more aligned matched-row activations than
        scrambled-row activations \emph{regardless} of any deeper
        substrate similarity.
    \item That the alignment is robust to held-out probes. We fit
        Procrustes on the full probe set and evaluate on the same set;
        we do not perform a train/test split.
    \item That the alignment exceeds the alignment achievable by simple
        lexical features (TF-IDF, sentence-transformer embeddings).
    \item That the alignment survives the shared training-distribution
        confound. Both models in each pair are trained on overlapping
        post-2024 web corpora.
    \item That the aligned subspace is causally responsible for any
        compositional behavior. Procrustes alignment is correspondence,
        not causation.
\end{enumerate}

We are explicit about each of these gaps because the dramatic
appearance of the headline numbers (residuals near $0.8$ against null
means near $1.0$, with empirical $p < 0.001$ at every cell) could be
read as much stronger evidence for the strong operator-theoretic claim
than the experiment supports. Section~\ref{sec:requiredwork} specifies
the additional studies that would, jointly, support the strong claim.
None of those studies has been run in the present paper.

\subsection*{Contributions}

\begin{enumerate}
    \item \textbf{An operator-theoretic framing as a research program.}
        We propose, and articulate at the level of the synthesis paper
        outlined elsewhere in this research program, that the
        mechanistic interpretability literature's empirical regularities
        are most parsimoniously read as evidence for approximately
        invariant subspaces of the forward operator. The framing is a
        program, not a finding.
    \item \textbf{Pre-registered methodology and infrastructure.} We
        specify a pre-registered protocol for cross-model
        representational alignment, with declared layer fractions,
        pooling, dimensionality threshold, significance criterion, and
        null distributions. The infrastructure is released for
        reuse and extension.
    \item \textbf{A representational-alignment pilot result.} Phases~0
        and 1 demonstrate that matched compositional-evaluation prompts
        produce more orthogonally alignable PCA-projected activation
        geometries across architecturally different frontier
        transformers than under random row permutation, with empirical
        $p < 0.001$ in every pre-registered cell.
    \item \textbf{An explicit accounting of what the pilot does not
        establish, and the studies that would close each gap.}
        Section~\ref{sec:requiredwork} catalogues five required
        additional studies (invariance-leakage measurement, lexical and
        template baselines, train/test Procrustes with cross-domain
        transfer, training-distribution-disjoint controls, causal
        cross-model intervention) and indicates which of them addresses
        which gap.
\end{enumerate}

\subsection*{Where This Sits in the Field}

We take the operator-theoretic framing seriously enough to claim it
unifies several open problems in interpretability that have been
discussed largely in isolation: the universality hypothesis
\cite{olah2020zoom, chughtai2023toy}, the polysemanticity / superposition
problem \cite{elhage2022superposition}, the cross-model
transfer-of-mechanism question implicit in the Platonic Representation
Hypothesis \cite{huh2024platonic}, and the substrate-independence
intuition that has organized the broader cognitive science of
representation since Marr \cite{marr1982vision}. We map these
explicitly in Section~\ref{sec:hardproblems}. The bridge to
computational neuroscience---to the neural-population doctrine
\cite{saxena2019towards} and to the geometry of biological invariants
\cite{okeefe1971hippocampus, hafting2005microstructure, gallego2017neural,
yamins2016using}---and to the cybernetic and free-energy frameworks
\cite{ashby1956introduction, friston2010free, achille2018emergence}
that anticipate the same structure from independent grounds, we touch
in the discussion (Section~\ref{sec:discussion}) but do not develop in
this paper. The full integration is the subject of a forthcoming
synthesis paper from the same research program.

\subsection*{Outline}

Section~\ref{sec:background} surveys the four literatures the bridge
draws on: operator theory, mechanistic interpretability, cross-model
representational alignment, and the substrate-independence tradition.
Section~\ref{sec:hardproblems} maps the open problems COSI sits
among. Section~\ref{sec:methods} formalizes the COSI hypothesis and
specifies the protocol. Section~\ref{sec:results} reports the Phase~0
calibration. Section~\ref{sec:limitations} catalogues the threats to
validity that Phase~0 does not address. Section~\ref{sec:discussion}
specifies Phases 1--3 and articulates the implications, conditional on
confirmation, across the literatures the framework bridges.
